%% ==========================================================================
%%  6  Goods first, and a reduction for doubly monotone valuations
%%
%%  New (2026-10-08). Universal chores can be added to any envy-free solution
%%  with unit subsidies at one more unit per agent. The procedure and the
%%  proofs mirror Sections 3.1 to 3.4 with every bundle replaced by its package
%%  (bundle plus subsidy). Gives a goods-first algorithm for objective signed
%%  valuations (two units per agent) and reduces doubly monotone dichotomous
%%  valuations to the insertion property (INS).
%% ==========================================================================
\section{Goods First, and a Reduction for Doubly Monotone Valuations}
\label{sec:goods-first}

Theorem~\ref{thm:objective-mixed} allocates the chores first and inserts the
goods afterwards. That order has no counterpart for doubly monotone valuations,
in which an item may be a good for some agents and a chore for others: such an
item cannot be set aside as a chore before it has been tried as a good. This
section reverses the order. We show that items that are chores for every agent
can be added at the end to any envy-free solution with unit subsidies, at the
cost of at most one more unit per agent and $n-1$ more in total
(Lemma~\ref{lem:universal-chores}). For objective signed dichotomous valuations
this gives a second, goods-first algorithm with at most two units per agent
(Theorem~\ref{thm:goods-first}). For doubly monotone dichotomous valuations it
reduces the same guarantee to a single insertion property for the remaining
items (Proposition~\ref{prop:dm-reduction}).

\begin{definition}[Doubly monotone dichotomous valuations, universal chores]
\label{def:dm}
A valuation profile $\set{v_i}_{i\in\agents}$ is \emph{doubly monotone
dichotomous} if for every agent $i$ there is a partition
$\items=G_i\mathbin{\dot\cup}C_i$ such that, for every $S\subseteq\items$,
\[
  v_i(S\cup\set g)-v_i(S)\in\set{0,1}
  \qquad\text{for every }g\in G_i\setminus S,
\]
and
\[
  v_i(S\cup\set c)-v_i(S)\in\set{0,-1}
  \qquad\text{for every }c\in C_i\setminus S.
\]
An item $e$ is a \emph{universal chore} if
$v_i(S\cup\set e)-v_i(S)\in\set{0,-1}$ for every agent $i$ and every
$S\subseteq\items\setminus\set e$.
\end{definition}

Objective signed dichotomous valuations (Definition~\ref{def:objective-signed})
are the case $G_i=G$ and $C_i=C$ for every $i$, and there every item of $C$ is a
universal chore. Doubly monotone dichotomous valuations are integer valued, by the
telescoping argument of Observation~\ref{obs:integrality}. Hence every envy graph
has integer arc weights, every pointwise-minimal subsidy vector is integral, and
Theorems~\ref{thm:hs-characterisation} and~\ref{thm:hs-minsubsidy} and
Proposition~\ref{prop:zero-coordinate} apply to allocations of any set of items,
since their proofs use only the arc weights.

% --------------------------------------------------------------------------
\subsection{Packages and the tie graph}
\label{sec:gf-ties}

Fix a set $M'\subseteq\items$, an allocation $A=(A_1,\dots,A_n)$ of $M'$ and an
integral subsidy vector $p$. Call $(A_j,p_j)$ the \emph{package} of agent $j$,
and write
\begin{equation}
\label{eq:gf-U}
  U_i(j):=v_i(A_j)+p_j
\end{equation}
for its value to agent $i$. Then $(A,p)$ is an envy-free solution if and only if
$U_i(i)\ge U_i(j)$ for all $i,j\in\agents$. All values $U_i(j)$ are integers, so
in an envy-free solution $U_i(j)\ne U_i(i)$ implies $U_i(j)\le U_i(i)-1$. The
\emph{tie graph} of $(A,p)$ is
\begin{equation}
\label{eq:gf-H}
  H=(\agents,E),
  \qquad
  (i,j)\in E \iff i\ne j\ \text{ and }\ U_i(i)=U_i(j).
\end{equation}
It is the equality graph \eqref{eq:m-graph} of \Cref{sec:m-terminal} with every
bundle replaced by its package. The payments must enter it. If agent $i$ prefers
$A_i$ to the bundle of a paid agent $j$ by exactly one, then $i$ is indifferent
between the two packages, and after one more costly chore $i$ would envy the
package of $j$. In the tie graph this indifference is the arc $(i,j)$.

% --------------------------------------------------------------------------
\subsection{Allocating the universal chores}
\label{sec:gf-procedure}

Let $(A,p)$ be an envy-free solution for an allocation $A$ of $M'\subseteq\items$
with $p\in\set{0,1}^n$, and suppose that every item of $R:=\items\setminus M'$ is
a universal chore. Starting from $(A,p)$, apply the following rules in the stated
order while $R\ne\emptyset$, recomputing $H$ after every step and removing from
$R$ every chore that is assigned. The procedure halts when $R=\emptyset$, or at
the final clause of (U3).
\begin{enumerate}
  \item[(U1)] If some $e\in R$ and some $i\in\agents$ satisfy
    $v_i(A_i\cup\set e)=v_i(A_i)$, assign $e$ to $i$.
  \item[(U2)] If some $e\in R$ and some arc $(i,j)\in E$ lying on a directed
    cycle of $H$ satisfy $v_i(A_j\cup\set e)=v_i(A_j)$, rotate the packages
    along that cycle, so that every agent on the cycle receives the bundle and
    the subsidy of the next agent on the cycle, and then assign $e$ to $i$.
  \item[(U3)] Otherwise select a strongly connected component $S$ of $H$ with no
    arc of $E$ leaving $S$. If $\abs R\ge\abs S$, assign one chore of $R$ to each
    agent of $S$, distinct agents receiving distinct chores. Otherwise halt.
\end{enumerate}
These are the rules (R1) to (R3) of \Cref{sec:m-terminal}, with the equality
graph replaced by the tie graph and with rotations that move each subsidy
together with its bundle.

\begin{lemma}
\label{lem:gf-step}
Each application of (U1), (U2) or (U3) assigns at least one chore of $R$ and
turns $(A,p)$ into an envy-free solution whose subsidy vector is a permutation of
the previous one. Consequently the procedure halts after at most
$\abs{\items\setminus M'}$ applications, and throughout it the subsidy vector
lies in $\set{0,1}^n$ and keeps its sum.
\end{lemma}

\begin{proof}
Every rule assigns at least one chore, (U3) assigning $\abs S\ge1$ of them, and
only (U2) moves subsidies, by a cyclic permutation. It remains to show that each
step preserves envy-freeness. Write $U'$ for the package values after the step.
Throughout we use that adding a universal chore to a bundle never raises any
agent's value for it.

\smallskip\noindent\textbf{(U1).} The package of $i$ keeps its value to $i$,
because $e$ is free for $i$ on $A_i$, and it is worth at most as much as before to
every other agent. All other packages are unchanged. Hence
$U'_k(k)=U_k(k)\ge U_k(j)\ge U'_k(j)$ for all $k,j\in\agents$.

\smallskip\noindent\textbf{(U2).} Let the cycle be
$i=a_0\to a_1\to\dots\to a_{L-1}\to a_0$, with $a_1=j$. After the step, an agent
$a_t$ with $t\ne0$ holds the package $(A_{a_{t+1}},p_{a_{t+1}})$, indices taken
modulo $L$, worth $U_{a_t}(a_{t+1})=U_{a_t}(a_t)$ to $a_t$ because
$(a_t,a_{t+1})\in E$. Agent $i$ holds $(A_j\cup\set e,p_j)$, worth
$v_i(A_j\cup\set e)+p_j=v_i(A_j)+p_j=U_i(j)=U_i(i)$ to $i$, because $e$ is free
for $i$ on $A_j$ and $(i,j)\in E$. Agents off the cycle keep their packages. So
every agent's own package keeps its value. The packages after the step are those
before it, except that $(A_j,p_j)$ has become $(A_j\cup\set e,p_j)$, which no
agent values more. Hence every agent $k$ values every package after the step at
most $\max_{j\in\agents}U_k(j)=U_k(k)=U'_k(k)$.

\smallskip\noindent\textbf{(U3).} Neither (U1) nor (U2) applies. Agents outside
$S$ keep their packages, and no package gains value to anyone, so every
$k\notin S$ satisfies $U'_k(k)=U_k(k)\ge U_k(j)\ge U'_k(j)$ for every $j$. Let
$i\in S$ and let $e_i$ be the chore assigned to $i$. As (U1) does not apply,
$e_i$ costs $i$ one unit on $A_i$, so $U'_i(i)=U_i(i)-1$. Let $j\ne i$.
\begin{itemize}
  \item If $j\notin S$, then $(i,j)\notin E$ because no arc leaves $S$. So
    $U_i(j)\le U_i(i)-1$, and $U'_i(j)=U_i(j)\le U'_i(i)$.
  \item If $j\in S$ and $(i,j)\in E$, then a path from $j$ to $i$ inside $S$
    closes a directed cycle through $(i,j)$, since $S$ is strongly connected. As
    (U2) does not apply, the chore $e_j$ assigned to $j$ costs $i$ one unit on
    $A_j$. So $U'_i(j)=U_i(j)-1=U_i(i)-1=U'_i(i)$.
  \item If $j\in S$ and $(i,j)\notin E$, then $U_i(j)\le U_i(i)-1$, and
    $U'_i(j)\le U_i(j)\le U'_i(i)$.
\end{itemize}
Thus the solution after the step is envy-free.
\end{proof}

\begin{lemma}
\label{lem:gf-terminal}
Suppose the procedure halts with $R\ne\emptyset$. Let $(A,p)$ be the solution at
which it halts, with tie graph $H=(\agents,E)$ and package values $U$, let
$r:=\abs R$, and let $S$ be the component selected by (U3) at the halting step.
Then the following hold.
\begin{enumerate}
  \item[(a)] $r<\abs S$.
  \item[(b)] $v_i(A_i\cup\set e)=v_i(A_i)-1$ for every $i\in\agents$ and every
    $e\in R$.
  \item[(c)] $v_i(A_j\cup\set e)=v_i(A_j)-1$ for every arc $(i,j)\in E$ with
    $i,j\in S$ and every $e\in R$.
  \item[(d)] $U_i(j)\le U_i(i)-1$ for every $i\in S$ and every
    $j\in\agents\setminus S$.
\end{enumerate}
\end{lemma}

\begin{proof}
The procedure halts at the final clause of (U3), so neither (U1) nor (U2) applies
and $\abs R<\abs S$, which is (a). Every chore of $R$ is universal, so its
marginals lie in $\set{0,-1}$. A marginal $0$ in (b) would let (U1) apply. In
(c), a path from $j$ to $i$ inside $S$ closes a directed cycle through $(i,j)$,
so a marginal $0$ would let (U2) apply. For (d), no arc leaves $S$, so
$U_i(j)\ne U_i(i)$, and envy-freeness gives $U_i(j)\le U_i(i)-1$.
\end{proof}

% --------------------------------------------------------------------------
\subsection{The completion}
\label{sec:gf-completion}

Suppose the procedure halts with $R=\set{e_1,\dots,e_r}\ne\emptyset$, and keep the
notation of Lemma~\ref{lem:gf-terminal}. By its part~(a) we may fix $r$ distinct
agents $T=\set{t_1,\dots,t_r}\subseteq S$, chosen arbitrarily. Define the complete
allocation $A^\ast$ by
\[
  A^\ast_{t_k}:=A_{t_k}\cup\set{e_k}\quad(1\le k\le r),
  \qquad
  A^\ast_i:=A_i\quad(i\notin T).
\]
Let $P\subseteq\agents$ be the smallest set such that
\begin{enumerate}
  \item[(P1)] $T\subseteq P$, and
  \item[(P2)] if $i\in\agents\setminus S$ and $(i,j)\in E$ for some $j\in P$,
    then $i\in P$,
\end{enumerate}
and set
\begin{equation}
\label{eq:gf-m}
  m_i:=p_i+\mathbf{1}[i\in P]
  \qquad (i\in\agents).
\end{equation}
These are the recipients and the subsidy set of \Cref{sec:m-completion}, built on
the tie graph. Every agent added to $P$ by (P2) lies outside $S$, so
$P\setminus T\subseteq\agents\setminus S$, exactly as in
Lemma~\ref{lem:m-Pstructure}.

\begin{lemma}
\label{lem:gf-expensive}
Let $i\in S$, let $t\in T$ and let $e\in R$. Then
$v_i(A_t\cup\set e)+p_t\le U_i(i)-1$.
\end{lemma}

\begin{proof}
Envy-freeness gives $U_i(t)\le U_i(i)$. If $U_i(t)<U_i(i)$, then
$U_i(t)\le U_i(i)-1$, and $v_i(A_t\cup\set e)\le v_i(A_t)$ because $e$ is a
universal chore. Otherwise $U_i(t)=U_i(i)$. If $i=t$, part~(b) of
Lemma~\ref{lem:gf-terminal} gives $v_i(A_i\cup\set e)+p_i=U_i(i)-1$. If $i\ne t$,
then $(i,t)\in E$ with $i,t\in S$, and part~(c) gives
$v_i(A_t\cup\set e)+p_t=U_i(t)-1=U_i(i)-1$.
\end{proof}

\begin{lemma}[Universal chores at one more unit per agent]
\label{lem:universal-chores}
Let $\set{v_i}_{i\in\agents}$ be doubly monotone dichotomous, let $A$ be an
allocation of a set $M'\subseteq\items$, and let $p\in\set{0,1}^n$ be such that
$(A,p)$ is an envy-free solution. Suppose that every item of
$\items\setminus M'$ is a universal chore. Then there are a complete allocation
$A^\ast$ and a subsidy vector $m\in\set{0,1,2}^n$ such that $(A^\ast,m)$ is an
envy-free solution and
\[
  \sum_{i\in\agents}m_i\le\sum_{i\in\agents}p_i+n-1 .
\]
More precisely, $m$ exceeds a permutation of $p$ by a vector in $\set{0,1}^n$.
Moreover, $(A^\ast,m)$ can be computed in polynomial time given value-oracle
access.
\end{lemma}

\begin{proof}
Run the procedure of \Cref{sec:gf-procedure}. By Lemma~\ref{lem:gf-step} it halts
at an envy-free solution whose subsidy vector, again written $p$, is a
permutation of the original one. If it halts with $R=\emptyset$, return this
solution with $m=p$. Otherwise construct $A^\ast$ and $m$ as in
\Cref{sec:gf-completion}. The allocation $A^\ast$ is complete, since the $t_k$
are distinct and every chore of $R$ is used once. The vector $m-p$ is the
indicator of $P$, so $m\in\set{0,1,2}^n$. Since
$P\subseteq T\cup(\agents\setminus S)$, $\abs T=r$ and $r\le\abs S-1$ by
Lemma~\ref{lem:gf-terminal}(a),
\[
  \sum_{i\in\agents}m_i=\sum_{i\in\agents}p_i+\abs P
  \le\sum_{i\in\agents}p_i+r+n-\abs S
  \le\sum_{i\in\agents}p_i+n-1 .
\]

It remains to show that $(A^\ast,m)$ is envy-free. Write
$U^\ast_i(j):=v_i(A^\ast_j)+m_j$. Since $A_j\subseteq A^\ast_j$ and the item
added to $A_j$, if any, is a universal chore,
\[
  U^\ast_i(j)\le U_i(j)+\mathbf{1}[j\in P],
  \qquad\text{with equality when }j\notin T.
\]
Fix $i,j\in\agents$. Since $T\subseteq P$, the following three cases are
exhaustive.

\smallskip\noindent\textbf{Case 1: $i\in T$.} Then $i\in S$ and $i\in P$, and
Lemma~\ref{lem:gf-terminal}(b) gives $U^\ast_i(i)=U_i(i)-1+1=U_i(i)$.
\begin{itemize}
  \item[(1a)] $j\in T$. Writing $e$ for the chore assigned to $j$,
    Lemma~\ref{lem:gf-expensive} gives
    $U^\ast_i(j)=v_i(A_j\cup\set e)+p_j+1\le U_i(i)$.
  \item[(1b)] $j\in P\setminus T$. Then $j\notin S$, so
    Lemma~\ref{lem:gf-terminal}(d) gives $U^\ast_i(j)=U_i(j)+1\le U_i(i)$.
  \item[(1c)] $j\notin P$. Then $U^\ast_i(j)=U_i(j)\le U_i(i)$.
\end{itemize}

\smallskip\noindent\textbf{Case 2: $i\in P\setminus T$.} Then $A^\ast_i=A_i$, so
$U^\ast_i(i)=U_i(i)+1\ge U_i(j)+1\ge U^\ast_i(j)$ for every $j$.

\smallskip\noindent\textbf{Case 3: $i\notin P$.} Then $i\notin T$, so
$U^\ast_i(i)=U_i(i)$.
\begin{itemize}
  \item[(3a)] $j\notin P$. Then $U^\ast_i(j)=U_i(j)\le U_i(i)$.
  \item[(3b)] $j\in P$ and $i\notin S$. If $U_i(j)=U_i(i)$, then $i\ne j$
    because $j\in P$ and $i\notin P$, so $(i,j)\in E$ and (P2) puts $i$ in $P$,
    a contradiction. Hence $U_i(j)\le U_i(i)-1$ and
    $U^\ast_i(j)\le U_i(j)+1\le U_i(i)$.
  \item[(3c)] $j\in P\setminus T$ and $i\in S$. Then $j\notin S$, so
    Lemma~\ref{lem:gf-terminal}(d) gives $U^\ast_i(j)=U_i(j)+1\le U_i(i)$.
  \item[(3d)] $j\in T$ and $i\in S$. Writing $e$ for the chore assigned to $j$,
    Lemma~\ref{lem:gf-expensive} gives
    $U^\ast_i(j)=v_i(A_j\cup\set e)+p_j+1\le U_i(i)$.
\end{itemize}
In every case $U^\ast_i(i)\ge U^\ast_i(j)$.

For the running time, the tie graph is determined by the $n^2$ values
$v_i(A_j)$, its strongly connected components and the cycles used by (U2) are
found by graph search, each rule tests at most $n^2\abs R$ marginals, and by
Lemma~\ref{lem:gf-step} there are at most $\abs{\items\setminus M'}$ rule
applications. Computing $P$ is a backward reachability search in $H$.
\end{proof}

% --------------------------------------------------------------------------
\subsection{A goods-first algorithm for objective signed valuations}
\label{sec:gf-objective}

\begin{theorem}
\label{thm:goods-first}
For every instance with objective signed dichotomous valuations, the following
procedure computes, in polynomial time given value-oracle access, a complete
allocation $\alloc$ and a subsidy vector $m\in\set{0,1,2}^n$ such that
$(\alloc,m)$ is envy-free and $\sum_{i\in\agents}m_i\le2(n-1)$. Allocate the
goods $G$ by the algorithm of Barman et al.~\cite{BKNS22}, replace its subsidy
vector by the pointwise-minimal one, and allocate the chores $C$ by
Lemma~\ref{lem:universal-chores}.
\end{theorem}

\begin{proof}
On subsets of $G$ every $v_i$ has marginals in $\set{0,1}$, so the algorithm
of~\cite{BKNS22} returns an allocation $A$ of $G$ and $p\in\set{0,1}^n$ with
$(A,p)$ envy-free. By Theorem~\ref{thm:hs-minsubsidy} the pointwise-minimal
subsidy vector is at most $p$ componentwise. It is integral, so it lies in
$\set{0,1}^n$, and by Proposition~\ref{prop:zero-coordinate} it has a zero
coordinate, so its sum is at most $n-1$. Every chore of $C$ is a universal chore,
so Lemma~\ref{lem:universal-chores} applies with $M'=G$ and gives
$m\in\set{0,1,2}^n$ with $\sum_i m_i\le(n-1)+(n-1)$. Both steps run in polynomial
time.
\end{proof}

Since the pointwise-minimal subsidy vector of $\alloc$ is at most $m$
componentwise, it satisfies the same two bounds. Theorem~\ref{thm:goods-first} is
weaker than Theorem~\ref{thm:objective-mixed}, which pays at most one unit per
agent. Its interest is the order of allocation, which is the order that
\Cref{sec:gf-dm} needs.

\begin{remark}[A chore with one unit of money attached]
\label{rem:gf-couples}
A tempting alternative completion attaches one unit of subsidy to each residual
chore and inserts the pair as a single item. The marginal value of the pair formed
from a chore $e$ is $v_i(S\cup\set e)-v_i(S)+1\in\set{0,1}$ for every agent and
every bundle, so the pair behaves as a dichotomous good. Formally, replace each
residual chore $e$ by a new item $\widehat e$, let $\widehat R$ be the set of new
items, and apply Lemma~\ref{lem:goods-extension} to the valuations
\[
  V_i(S):=v_i\bigl((S\setminus\widehat R)\cup\set{e:\widehat e\in S}\bigr)
  +\abs{S\cap\widehat R}.
\]
This variant keeps the total bound. At most $n-1$ pairs are formed, since the
procedure of \Cref{sec:gf-procedure} halts with $r<\abs S\le n$, and
Lemma~\ref{lem:goods-extension} keeps the remaining subsidy vector in
$\set{0,1}^n$, with a zero coordinate after passing to the pointwise-minimal
vector. So the payments total at most $2(n-1)$. It loses the per-agent bound,
because the units attached to the pairs accumulate when one agent receives
several of them, as Example~\ref{ex:gf-couples} shows. The distinct recipients
of \Cref{sec:gf-completion} prevent exactly this.
\end{remark}

\begin{example}
\label{ex:gf-couples}
Take four agents $x,a,b,d$, goods $g_1,g_2,g_3$ and chores $c_1,c_2$, and give
every agent the additive valuation $v(S)=\abs{S\cap G}-\abs{S\cap C}$. In the
algorithm of~\cite{BKNS22} every good is extendable, so it goes to a maximally
subsidized agent with marginal value one, and the maximally subsidized agents are
exactly those holding no good. So, up to renaming, the goods phase returns
$A_x=\emptyset$, $A_a=\set{g_1}$, $A_b=\set{g_2}$ and $A_d=\set{g_3}$, with
subsidy $1$ for $x$ and $0$ for the others. Every agent values every package at
$1$, so the tie graph is complete and $S=\agents$. Since $\abs R=2<4$ and every
chore costs every agent one unit on every bundle, the procedure of
\Cref{sec:gf-procedure} halts at once.

Each pair $\widehat c_k$ has marginal value $0$ for every agent on every bundle,
so it is not extendable, and \textsc{FindSink} starts at $x$, the only maximally
subsidized agent. Giving $\widehat c_1$ to $x$ changes no agent's value for any
package, so the minimal subsidies stay $(1,0,0,0)$, no agent reaches $2$, and
\textsc{FindSink} returns $x$. The same happens for $\widehat c_2$. Agent $x$
ends with $\set{c_1,c_2}$, two attached units and a subsidy of $1$, so $x$ is paid
$3$. The final allocation forces this, since $v_x(A_a)-v_x(\set{c_1,c_2})=3$, so
by Theorem~\ref{thm:hs-minsubsidy} no envy-free subsidy vector pays $x$ less than
$3$. With $n$ agents, $n-1$ such goods and $n-2$ such chores, the same run pays
$x$ exactly $n-1$.

The completion of \Cref{sec:gf-completion} instead gives $c_1$ and $c_2$ to two
distinct agents of $S$. For $T=\set{a,b}$ it pays $x$, $a$ and $b$ one unit each,
and every agent ends with value plus subsidy equal to $1$.
\end{example}

% --------------------------------------------------------------------------
\subsection{Doubly monotone valuations}
\label{sec:gf-dm}

Lemma~\ref{lem:universal-chores} needs nothing from the items allocated before it
beyond an envy-free solution with unit subsidies. This reduces the two-unit
guarantee for doubly monotone dichotomous valuations to the following insertion
property.
\begin{enumerate}
  \item[(INS)] For every set $M'\subseteq\items$, every allocation $A$ of $M'$
    and every $p\in\set{0,1}^n$ such that $(A,p)$ is envy-free, and every item
    $e\in\items\setminus M'$ that is not a universal chore, there exist an
    allocation $A'$ of $M'\cup\set e$ and $p'\in\set{0,1}^n$ such that
    $(A',p')$ is envy-free.
\end{enumerate}

\begin{proposition}
\label{prop:dm-reduction}
Let $\set{v_i}_{i\in\agents}$ be doubly monotone dichotomous and satisfy (INS).
Then there are a complete allocation $\alloc$ and a subsidy vector
$m\in\set{0,1,2}^n$ such that $(\alloc,m)$ is envy-free and
$\sum_{i\in\agents}m_i\le2(n-1)$.
\end{proposition}

\begin{proof}
Start from the empty allocation with zero subsidies, which is envy-free. Insert
the items that are not universal chores one at a time by (INS), each time
replacing the subsidy vector by the pointwise-minimal one. That vector is
integral and componentwise at most the vector given by (INS), so it lies in
$\set{0,1}^n$. When these items are exhausted, the current solution is envy-free,
its subsidy vector lies in $\set{0,1}^n$ and has sum at most $n-1$ by
Proposition~\ref{prop:zero-coordinate}, and every unallocated item is a universal
chore. Lemma~\ref{lem:universal-chores} completes it with $m\in\set{0,1,2}^n$ and
$\sum_i m_i\le(n-1)+(n-1)$.
\end{proof}

Lemma~\ref{lem:goods-extension} proves (INS) for every item that is a good for
every agent, that is, whose marginal value lies in $\set{0,1}$ for every agent and
every bundle, whatever the current bundles contain. For objective signed
valuations every item that is not a universal chore is of this kind, so
Proposition~\ref{prop:dm-reduction} also yields the bound of
Theorem~\ref{thm:goods-first}. The open case is an item that is a good for some
agents and a chore for others. The argument behind
Lemma~\ref{lem:goods-extension} uses at several points that the new item has a
nonnegative marginal value for the agent receiving it, for instance in Claim~2
and in the inequality $\operatorname{SW}(X^{s_0})\geq\operatorname{SW}(A)$ of
Case~2 in Appendix~\ref{app:goods-extension}, and this fails when such an item is
tentatively given to an agent for whom it is a chore.
Whether (INS) holds for such items is open.
